% Generated draft: requires mathematical and source review.
\documentclass[11pt]{article}
\usepackage[margin=25mm]{geometry}
\usepackage{fontspec}
\setmainfont{Latin Modern Roman}
\usepackage{amsmath,amssymb,mathtools,unicode-math}
\setmathfont{Latin Modern Math}
\usepackage{enumitem,needspace,xurl,hyperref}
\setlength{\parindent}{0pt}
\setlength{\parskip}{5pt}
\setlength{\emergencystretch}{3em}
\newcommand{\N}{\mathbb{N}}
\newcommand{\Z}{\mathbb{Z}}
\newcommand{\Q}{\mathbb{Q}}
\newcommand{\R}{\mathbb{R}}
\newcommand{\C}{\mathbb{C}}
\newcommand{\F}{\mathbb{F}}
\newcommand{\A}{\mathbb{A}}
\newcommand{\PP}{\mathbb{P}}
\newcommand{\E}{\mathbb{E}}
\newcommand{\eps}{\varepsilon}
\newcommand{\dd}{\,\mathrm{d}}
\begin{document}
\section*{UnsolvedMath ID: 30006547}
\textbf{Source problem code:} catalogue entry
\subsection{Definitions and mathematical statement}
Let \(r\) be a positive integer, and write \([r]=\{1,\ldots,r\}\). An \(r\)-edge-colored graph is a finite simple graph \(G=(V,E)\) together with a map \(\chi:E\to[r]\). Here the coloring is not required to be proper or to use every color; these conventions are not explicitly specified in the supplied record. A rainbow triangle is a three-element set \(\{u,v,w\}\subseteq V\) such that all three edges \(uv,vw,wu\) belong to \(E\) and their colors are pairwise distinct. Let \(T(G,\chi)\) denote the number of such sets.

For each positive integer \(m\), define
\[
M_r(m)=\max\bigl\{T(G,\chi): |E(G)|=m,\ \chi:E(G)\to[r]\bigr\}.
\]
One precise interpretation of the requested sharp constant is the optimal universal coefficient
\[
C_r=\sup_{m\geq1}\frac{M_r(m)}{m^{3/2}},
\]
so that the intended inequality is
\[
T(G,\chi)\leq C_r\,|E(G)|^{3/2}
\]
for every such colored graph with at least one edge, with \(C_r\) as small as possible. Graphs with no edges have no triangles. The supplied record does not explicitly distinguish this universal coefficient from an asymptotic leading coefficient.

For an integer \(q\geq2\), an affine plane of order \(q\) is an incidence structure with \(q^2\) points whose lines are partitioned into \(q+1\) parallel classes, each consisting of \(q\) disjoint lines covering all points, with \(q\) points on every line and exactly one line through any two distinct points.

\textbf{Problem.} Determine \(C_r\) for each integer \(r\geq3\) for which no affine plane of order \(r-1\) exists. In particular, determine \(C_7\) and \(C_{11}\), corresponding to orders \(6\) and \(10\).

The range \(r\geq3\) uses the conventional nondegenerate affine-plane orders \(q\geq2\); the supplied statement does not state a lower bound on \(r\). Its original excerpt contains a damaged upper-bound formula, so no numerical general upper bound is asserted here. This is a statement-only interpretation requiring mathematical review.
\subsection{Short English statement}
Count triangles whose three edges have distinct colors in a simple graph colored from a palette of \(r\) colors. Determine the smallest universal coefficient multiplying \(m^{3/2}\) that bounds this count when no affine plane of order \(r-1\) exists, especially for \(r=7\) and \(r=11\). The universal-coefficient interpretation and coloring conventions require confirmation against the source.
\subsection{Sources}
\begin{itemize}
\item Local UnsolvedMath registry, permanent ID \(30006547\), title: \emph{Sharp Constants for Multicolor Rainbow Triangles}; supplied complete source record, source code \texttt{catalogue entry}, in the Oberwolfach Reports Open Problems collection. No individual proposer is supplied; the proposed year is recorded as 2026.
\item \emph{Combinatorics} (2026), Oberwolfach Reports, DOI 10.4171/owr/2026/1, pp.\ 5--79. Supplied source citation and literature reference labeled ``Oberwolfach Report 1/2026'': \url{https://doi.org/10.4171/owr/2026/1}. The supplied background describes these reports as records of mathematical workshops and their open-problem discussions. No problem-specific page is supplied.
\item Chao--Yu, \emph{A purely entropic approach}: \url{https://arxiv.org/abs/2407.14084}. This reference is supplied in the literature sources; its claims and relevance have not been independently checked here.
\item Literature assessment dated \texttt{2026-08-22}: supplied dated curation, not independently verified. It labels the problem open and states that the general upper bound and prime-power constructions do not determine the exact constant when the required affine plane does not exist. The supplied background requests independent verification of time-sensitive status and cited claims.
\item The supplied record labels its statement \texttt{corrected\_verified} and reports that the official report was checked, the question isolated, and its notation repaired. These are supplied verification claims, not verification performed for this draft. The original excerpt attributes a mixture bound used in a proof to Maya Sankar; it does not identify her as the proposer of this problem.
\end{itemize}
\end{document}
